\chapter{\texorpdfstring{\latestrevision{Complete-Policy Transfer and Target Adaptation}}{Complete-Policy Transfer and Target Adaptation}}
\label{ch:screen-f-transfer}

\latestrevision{Chapter~\ref{ch:encoder-transfer-bridge} tested whether a source
encoder remains useful after its incompatible action heads are replaced. This
chapter asks the stricter question: can the shared candidate-action actor
learned on \texttt{bus14} control WCCI while retaining both its encoder and its
action-selection interface? It first evaluates the complete actor without
target training, then studies decision-time changes and target adaptation.}

\section{Comparison rules}
\label{sec:wcci-design}

\latestrevision{The transferred object is the complete candidate-action actor,
not only its graph encoder.} The actor varies along four binary choices: original or corrected inputs, mean
or typed-mean candidate pooling, action descriptor off or on, and a shared or
dedicated idle-action head. Table~\ref{tab:wcci-transfer-design} defines the
three training routes.

\begin{table}[H]
  \centering
  \small
  \caption{Training routes used in the transfer study.}
  \label{tab:wcci-transfer-design}
  \setlength{\tabcolsep}{5pt}
  \begin{tabular}{L{0.22\textwidth}L{0.30\textwidth}L{0.38\textwidth}}
    \toprule
    Route & Initial actor & WCCI training \\
    \midrule
    Zero-shot & Shared \texttt{bus14} actor & None \\
    From scratch & Random weights & 15M target steps \\
    Fine-tuned & Shared \texttt{bus14} actor & About 5M target steps \\
    \bottomrule
  \end{tabular}
\end{table}

Final checkpoints are used for actors trained on WCCI. Zero-shot evaluation
uses the source checkpoint selected on \texttt{bus14} \revision{by mean survival
over the 20 training-split chronics in each periodic rollout, with the chronic
set changing between evaluations}. The optional heuristic
blocks an agent's action when its local maximum line loading is below $0.95$.
Every comparison uses the same 50 chronics and seed 0.

\section{What controls zero-shot transfer?}
\label{sec:wcci-zero-shot-determinants}

Figure~\ref{fig:wcci-zero-shot-determinants} separates three questions. The top
panel gives mean survival over all 50 chronics. The lower-left panel restricts
the score to the 24 difficult chronics. The lower-right panel then adjusts that
difficult-cohort score to the greedy reference available at the same $k$:
\begin{equation}
  R_{k,\pi}=100\,
  \frac{S_{\pi}-S_{\mathrm{DN}}}
       {S_{\mathrm{greedy},k}-S_{\mathrm{DN}}}.
  \label{eq:wcci-greedy-progress}
\end{equation}
Here $0\%$ means do-nothing performance and $100\%$ means the greedy result at
the same action count. A negative value is worse than do nothing. This avoids
crediting a policy simply because it receives a larger candidate list.

\begin{figure}[H]
  \centering
  \includegraphics[width=\linewidth]{figures/wcci_zero_shot_determinants.png}
  \caption{Absolute and feasible-relative zero-shot performance. Each box
  contains all 16 NL/NLS architectures. Green stars give the greedy full-test
  score and dashed lines give do nothing. The lower-right score measures
  progress from do nothing to the greedy reference at the same action-space
  size.}
  \label{fig:wcci-zero-shot-determinants}
\end{figure}

The distributions show that absolute and feasible-relative performance answer
different questions. With the heuristic, the \revision{median of the models'
full-test mean-survival scores} is $62.5\%$ at $k=32$ and about $54.5\%$ for the
larger action spaces; the \revision{median difficult-cohort mean survival across
models} falls from $21.9\%$ to roughly $11\%$. Relative to the
same-$k$ greedy reference, the median recovers $54.6\%$ of the available gain at
$k=32$, but only $5.8\%$, $5.2\%$ and $2.9\%$ at $k=64,128,256$. The gate helps,
but the actors exploit a smaller fraction of what becomes feasible as the action
set grows. Ungated median performance remains below do nothing at every $k$.

Figure~\ref{fig:wcci-zero-shot-architecture} then separates the architecture
effects by gate condition. Typed-mean \revision{candidate pooling} is the clearest positive effect:
$+5.4$ points ungated and $+8.7$ points gated. Physical scaling adds about one
point in both cases. The descriptor is small on average, while the dedicated
idle head is negative ungated and nearly neutral gated.

\begin{figure}[H]
  \centering
  \includegraphics[width=\linewidth]{figures/wcci_zero_shot_architecture_effects.png}
  \caption{Matched architecture effects on the difficult cohort, shown
  separately for ungated actors and the local-$\rho=0.95$ heuristic. Positive
  values improve survival. Error bars are bootstrap intervals over matched
  architecture cells; $n$ is the number of available pairs.}
  \label{fig:wcci-zero-shot-architecture}
\end{figure}

These comparisons use one seed. They identify useful directions but do not
establish a final architecture ranking.

\section{Policy changes and the local heuristic}

Two changes are tested around the zero-shot actor. The first uses the
global-maximum readout and exact action-delta encoder of
Figure~\ref{fig:candidate-action-gmax-delta-pipeline}. The second uses an
adaptive intervention budget. Figure~\ref{fig:wcci-policy-change-performance}
first shows their absolute performance in one view. Each action-space cluster
contains four distributions: the two policy changes, with and without the
local-$\rho$ heuristic.

\begin{figure}[H]
  \centering
  \includegraphics[width=\linewidth]{figures/wcci_policy_change_absolute_performance.png}
  \caption{Absolute performance of the Gmax-delta and adaptive-budget actors
  at each action-space size, with four distributions per cluster. Boxes and
  points summarize the architecture cells for ungated and local-$\rho=0.95$
  evaluation. The dashed line marks do-nothing; the green reference gives the
  greedy ceiling or, in the relative panel, its full attainable gain.}
  \label{fig:wcci-policy-change-performance}
\end{figure}

The absolute view shows whether a positive change is large enough to produce a
useful policy. On difficult chronics, ungated Gmax-delta falls from a $19.1\%$
median at $k=32$ to $5.1\%$ at $k=256$, whereas the adaptive budget retains
$12.7\%$ at $k=256$. With the heuristic, both modified actors remain above the
do-nothing reference at every $k$. Figure~\ref{fig:wcci-intervention-gate} then
isolates the matched change from the plain actor and the effect of the
local-$\rho$ heuristic.

\begin{figure}[H]
  \centering
  \includegraphics[width=\linewidth]{figures/wcci_intervention_gate_effects.png}
  \caption{Matched changes on the difficult cohort. Positive values are better.
  The Gmax-delta and adaptive-budget blocks are compared with the same plain
  zero-shot architectures.}
  \label{fig:wcci-intervention-gate}
\end{figure}

Without the heuristic, Gmax-delta adds $4.6$ points and the adaptive budget
adds $6.0$ points. Across gate conditions, the heuristic adds $7.2$ points to
the plain zero-shot actor and $5.3$ points to both modified actors. It is not a
universal safety layer: it reduces the scratch scorer by $2.2$ points and the
flat MLP by up to $2.8$ points. The gate must therefore be checked with each
policy.

The plain block contains all 64 matched architecture--$k$ pairs. Gmax-delta and
the adaptive budget each contain 60 pairs because their complete designs use 15
architectures rather than 16.

\section{Does candidate-policy target training help?}

Figure~\ref{fig:wcci-adaptation} compares the same eight NLS architectures at
$k=64$. The left panel tracks each architecture across zero-shot, scratch and
fine-tuned routes. The right panel shows the main trade-off: \revision{mean survival} on the
difficult cohort against the number of easy chronics kept.

\begin{figure}[H]
  \centering
  \includegraphics[width=\linewidth]{figures/wcci_adaptation_summary.png}
  \caption{Effect of target-grid training. Grey lines join the same architecture
  in the left panel. The black cross and green star in the right panel are do
  nothing and the greedy $k=64$ reference.}
  \label{fig:wcci-adaptation}
\end{figure}

Fine-tuning is the most consistent route. It beats training from scratch on
seven of the eight difficult-cohort comparisons, with a mean gain of $7.9$
points. It also beats zero-shot transfer in all eight cells, by $12.0$ points on
average. The \revision{median difficult-cohort mean survival across architectures} rises from $5.7\%$ zero-shot to
$9.6\%$ from scratch and $18.1\%$ after fine-tuning.

The right panel adds an important qualification. Some actors improve difficult
chronics while losing cases that do nothing already solves. The useful region
is therefore the upper-right corner, not simply the highest difficult-cohort
score. The target-trained controls remain far below the greedy policy, showing
that much more performance is available in the action space.

\section{Limits of the study}

The study uses one training seed and 50 evaluation chronics. Bootstrap intervals
describe variation across architecture cells, not across independent training
runs. Some campaigns also use different target-training budgets. The results
support comparisons inside the matched blocks, but small gaps between unrelated
blocks should not be overinterpreted. Full matrices and exact control tables are
provided in Appendix~\ref{app:wcci-result-atlas}.

\begin{tcolorbox}[
  colframe=red!75!black,
  colback=red!5!white,
  title=Chapter 10 Summary: Complete-Policy Transfer and Adaptation]
\textbf{Purpose:} \latestrevision{To determine how well the complete shared
candidate-action actor learned on \texttt{bus14} transfers to WCCI, identify
which architectural and decision-time choices improve zero-shot behavior, and
measure whether WCCI training benefits more from source initialization than
from random weights.}
\\
\\
\textbf{Content:} \latestrevision{Unlike Chapter~\ref{ch:encoder-transfer-bridge},
the source action interface is not replaced.} We evaluate 16 NL/NLS architecture cells over
$k\in\{32,64,128,256\}$, with and without a local-$\rho=0.95$ gate, using both
absolute \revision{mean survival} and progress from do nothing to the same-$k$ greedy reference.
We then test Gmax-delta and adaptive-budget policy changes before comparing the
same eight NLS architectures at $k=64$ under zero-shot transfer, training from
scratch and target fine-tuning.
\\
\\
\textbf{Findings:}
\begin{itemize}
  \item Zero-shot transfer is strongest with the smallest action set. With the
  gate, $k=32$ reaches $62.5\%$ overall and $21.9\%$ on difficult chronics,
  recovering $54.6\%$ of the available same-$k$ gain; the recovered fraction
  falls to $2.9\%$ at $k=256$. Ungated medians remain below do nothing.
  \item Typed-mean \revision{candidate pooling} is the clearest positive architecture effect,
  adding $5.4$ \revision{difficult-cohort mean-survival points} ungated and $8.7$ gated. Physical
  scaling adds about one point, while the descriptor is small on average and a
  dedicated idle head is harmful ungated and nearly neutral gated.
  \item Gmax-delta and the adaptive budget add $4.6$ and $6.0$ matched points
  without the gate. The gate improves 58 of 64 plain cells and adds $7.2$ points
  on average, but it hurts the scratch scorer and flat MLP; it is therefore
  policy-dependent rather than a universal safety mechanism.
  \item Fine-tuning is the strongest adaptation route. It beats zero-shot in
  all eight matched cells by $12.0$ points on average and beats scratch in seven
  cells by $7.9$ points. \revision{Median difficult-cohort mean survival across architectures} progresses from
  $5.7\%$ zero-shot to $9.6\%$ scratch and $18.1\%$ fine-tuned.
\end{itemize}
\textbf{Takeaway:} Reliable transfer requires controlling the action-selection
problem, not merely enlarging the available action set. Small candidate sets
and typed-mean \revision{candidate pooling} are the most robust zero-shot choices; local gating must
be validated per policy; and conservative target fine-tuning is more effective
than either zero-shot deployment or training the same actor from scratch. These
conclusions are based on one seed and 50 chronics, so matched effects are more
credible than small gaps between unrelated cells.
\end{tcolorbox}
